\documentclass[11pt]{article}
\usepackage[localtheorems]{raeez-math-template}
\newtheorem{theorem}{Theorem}[section]
\newtheorem{proposition}[theorem]{Proposition}
\newtheorem{lemma}[theorem]{Lemma}
\newtheorem{definition}[theorem]{Definition}
\newtheorem{remark}[theorem]{Remark}
\title{The Shadow Tower: $A_\infty$ Coproduct Corrections and Feynman Coefficients}
\author{Raeez Lorgat}
\date{2026-04-22}
\begin{document}
\maketitle

\begin{abstract}
The stress-tensor line in the $c=1$ Virasoro subalgebra of the
$\mathcal W_{1+\infty}$ completion of the affine CY$_3$ Hall theory has
a single quadratic shadow landscape. With
$\kappa_{\mathrm{ch}}=c/2$, cubic coefficient $\alpha=2c$, and
$S_4=10/(c(5c+22))$, the positive branch
\[
  f(t)^2
  =
  4\kappa_{\mathrm{ch}}^2
  +12\kappa_{\mathrm{ch}}\alpha t
  +(9\alpha^2+16\kappa_{\mathrm{ch}}S_4)t^2
\]
determines all post-quartic $A_\infty$ coefficients. At $c=1$ this
gives
\[
  m_3=-2,\quad m_4=\frac{40}{27},\quad
  m_5=-\frac{80}{9},\quad m_6=\frac{38080}{729},\quad
  m_7=-\frac{24320}{81},\quad
  m_8=\frac{33157760}{19683},
\]
and
\[
  S_5=-\frac{16}{9},\quad S_6=\frac{19040}{2187},\quad
  S_7=-\frac{24320}{567},\quad
  S_8=\frac{4144720}{19683}.
\]
The connected Wick cumulant
$G_5^{\mathrm{conn}}(1,2,3,4,5)=775/5184$ is an independent finite
enumeration of ten half-edge pairings. It witnesses the same Virasoro
finite-OPE input, while the conversion to $m_5$ passes through the
chosen residue and bar-normalisation map. The Hall-side equality is
$\mathrm{CoHA}(\mathbb C^3)=Y^+(\widehat{\mathfrak{gl}}_1)$; the
$\mathcal W_{1+\infty}$ line belongs to the completed double or centre
layer.
\end{abstract}

\section{The $c=1$ Virasoro Line}

The affine CY$_3$ specialisation of $\Phi$ supplies an $E_1$ chiral
Hall algebra. Its positive Hall half is
\[
  \mathrm{CoHA}(\mathbb C^3)=Y^+(\widehat{\mathfrak{gl}}_1).
\]
The vertex algebra $\mathcal W_{1+\infty}$ appears after passing to the
completed Drinfeld double, equivalently to the centre or representation
category layer. This separation fixes the operadic scope: for $d\geq3$
the algebra $A$ is $E_1$, and the braided $E_2$ structure is carried by
$Z(\mathrm{Rep}(A))$ or by the corresponding completed double.

The self-dual Omega point
\[
  h_1=1,\qquad h_2=0,\qquad h_3=-1
\]
contains the Virasoro stress-tensor line
$L_T=\mathbf{k}\cdot T$ at central charge $c=1$. All scalar statements
below are made on this one-dimensional line and use the normalisation
\[
  \langle T(z)T(w)\rangle=\frac{c/2}{(z-w)^4},\qquad
  \kappa_{\mathrm{ch}}=\frac c2.
\]

\section{OPE Inputs}

\begin{definition}[The three scalar inputs]
\label{def:three-inputs}
For the Virasoro stress tensor at central charge $c$, set
\[
  \kappa_{\mathrm{ch}}=\frac c2,\qquad
  \alpha=2c,\qquad
  S_4=\frac{10}{c(5c+22)}.
\]
The coefficient $\alpha$ is read from the associator convention
\[
  m_3(T,T,T)=-\,\alpha T.
\]
\end{definition}

\begin{proposition}[Level-four denominator]
\label{prop:level-four-denominator}
The quartic shadow input is
\[
  S_4=\frac{10}{c(5c+22)}.
\]
\end{proposition}

\begin{proof}
In the Virasoro vacuum module at level $4$, use the basis
$L_{-4}|0\rangle$, $L_{-2}^2|0\rangle$. The commutation relation
\[
  [L_m,L_n]=(m-n)L_{m+n}+\frac{c}{12}(m^3-m)\delta_{m+n,0}
\]
gives the Gram matrix
\[
  G_4=
  \begin{pmatrix}
    5c & 3c\\
    3c & c(c+8)/2
  \end{pmatrix}.
\]
The quasi-primary
\[
  \Lambda=L_{-2}^2|0\rangle-\frac35 L_{-4}|0\rangle
\]
has norm
\[
  \langle\Lambda,\Lambda\rangle
  =\frac{c(c+8)}2-\frac{(3c)^2}{5c}
  =\frac{c(5c+22)}{10}.
\]
The inverse norm is the Zamolodchikov quartic channel coefficient, so
$S_4=10/(c(5c+22))$.
\end{proof}

At $c=1$,
\[
  \kappa_{\mathrm{ch}}=\frac12,\qquad
  \alpha=2,\qquad
  S_4=\frac{10}{27}.
\]

\section{Newton Shadow Recursion}

\begin{definition}[Shadow landscape and coefficients]
\label{def:shadow-landscape}
Let
\[
  Q_c(t)
  =
  4\kappa_{\mathrm{ch}}^2
  +12\kappa_{\mathrm{ch}}\alpha t
  +(9\alpha^2+16\kappa_{\mathrm{ch}}S_4)t^2,
\]
and let
\[
  f(t)=\sqrt{Q_c(t)}=\sum_{n\geq0}a_nt^n
\]
be the branch with $a_0=2\kappa_{\mathrm{ch}}$. For $k\geq5$ the
post-quartic $A_\infty$ coefficient and the shadow coefficient are
\[
  m_k(T,\ldots,T)=a_{k-2}T,\qquad
  S_k=\frac{a_{k-2}}{k}.
\]
The lower arities are the OPE normalisations
\[
  m_3=-\alpha,\qquad S_3=\alpha,\qquad m_4=4S_4.
\]
\end{definition}

\begin{theorem}[Square-root recurrence]
\label{thm:square-root-recurrence}
The coefficients of $f(t)$ satisfy
\[
  a_0=2\kappa_{\mathrm{ch}},\qquad
  a_1=\frac{12\kappa_{\mathrm{ch}}\alpha}{2a_0},\qquad
  a_2=\frac{9\alpha^2+16\kappa_{\mathrm{ch}}S_4-a_1^2}{2a_0},
\]
and, for $n\geq3$,
\[
  a_n=-\frac{1}{2a_0}\sum_{i=1}^{n-1}a_i a_{n-i}.
\]
\end{theorem}

\begin{proof}
Expand $f(t)^2=Q_c(t)$. The coefficients of $t^0$, $t^1$, and $t^2$
give the first three displayed equations. Since $Q_c(t)$ is quadratic,
the coefficient of $t^n$ vanishes for every $n\geq3$, giving
\[
  2a_0a_n+\sum_{i=1}^{n-1}a_i a_{n-i}=0.
\]
\end{proof}

\begin{proposition}[Equivalent $S_r$ recurrence]
\label{prop:sr-recurrence}
For $r\geq5$,
\[
  S_r
  =
  -\frac{1}{2r\kappa_{\mathrm{ch}}}
  \sum_{\substack{j+k=r+2\\3\leq j\leq k<r}}
  \epsilon(j,k)\,jk\,S_jS_k,
\]
where $\epsilon(j,k)=1$ for $j<k$ and $\epsilon(j,j)=1/2$.
\end{proposition}

\begin{proof}
Substitute $a_{r-2}=rS_r$ into
Theorem~\ref{thm:square-root-recurrence}. The restriction
$j\leq k$ records unordered pairs; the factor $\epsilon(j,j)=1/2$
removes the double count on the diagonal.
\end{proof}

\section{Exact Tower Through Arity Eight}

\begin{theorem}[The $c=1$ shadow tower through $m_8$]
\label{thm:c1-shadow-tower-through-eight}
At $c=1$,
\[
\begin{array}{c|c|c}
k & m_k(T,\ldots,T)/T & S_k\\
\hline
3 & -2 & 2\\
4 & 40/27 & 10/27\\
5 & -80/9 & -16/9\\
6 & 38080/729 & 19040/2187\\
7 & -24320/81 & -24320/567\\
8 & 33157760/19683 & 4144720/19683
\end{array}
\]
\end{theorem}

\begin{proof}
At $c=1$,
\[
  Q_1(t)=1+12t+\frac{1052}{27}t^2.
\]
Thus $a_0=1$, $a_1=6$, and
\[
  a_2=\frac{1052/27-36}{2}=\frac{40}{27}.
\]
The recurrence gives
\[
  a_3=-\frac{2\cdot6\cdot(40/27)}2=-\frac{80}{9},
\]
\[
  a_4=-\frac{2\cdot6\cdot(-80/9)+(40/27)^2}{2}
      =\frac{38080}{729},
\]
\[
  a_5=-\frac{2\cdot6\cdot(38080/729)
        +2\cdot(40/27)\cdot(-80/9)}2
      =-\frac{24320}{81},
\]
and
\[
  a_6=-\frac{2\cdot6\cdot(-24320/81)
        +2\cdot(40/27)\cdot(38080/729)
        +(-80/9)^2}{2}
      =\frac{33157760}{19683}.
\]
Definition~\ref{def:shadow-landscape} converts $a_{k-2}$ to $m_k$ and
$S_k$ for $k\geq5$; the $k=3,4$ rows are the OPE normalisations.
\end{proof}

\begin{remark}[Denominators]
\label{rem:denominators}
The post-quartic operation coefficients $m_k/T=a_{k-2}$ displayed
above have denominators that are powers of $3$. The only denominator in
the $c=1$ input data is $5c+22=27$ and the recurrence divides by
$2a_0=2$ after a symmetric convolution whose numerator is even. The
shadow coefficients $S_k=a_{k-2}/k$ may acquire an arity denominator
before cancellation; $S_7=-24320/567$ retains the factor $7$.
\end{remark}

\section{Convergence Radius}

\begin{proposition}[Branch points at $c=1$]
\label{prop:convergence-radius}
The positive branch of
\[
  f(t)=\sqrt{1+12t+\frac{1052}{27}t^2}
\]
has convergence radius
\[
  \rho=\sqrt{\frac{27}{1052}}
       =\frac{3\sqrt3}{2\sqrt{263}}
       \approx 0.1602042423.
\]
\end{proposition}

\begin{proof}
The branch points are the two roots of
$1+12t+(1052/27)t^2$. Their product is $27/1052$. The discriminant is
\[
  12^2-4\cdot\frac{1052}{27}
  =-\frac{320}{27},
\]
so the two roots are complex conjugates and have common modulus
$\sqrt{27/1052}$.
\end{proof}

The tower is algebraic and non-terminating, hence class $\mathbf M$ on
this line. The exponential scale is $\rho^{-1}\approx6.242031958$,
with the usual algebraic square-root subexponential factor.

\section{Connected Five-Point Wick Cumulant}

\begin{definition}[Exact connected Wick sum]
\label{def:exact-wick-sum}
At $c=1$ write
\[
  T(z)=-\frac12{:}J(z)^2{:},\qquad
  \langle J(z)J(w)\rangle=-\frac1{(z-w)^2}.
\]
For $k$ marked points, attach two labelled half-edges to each point.
Let $\mathcal P_k^{\mathrm{conn}}$ be the set of perfect matchings of
these $2k$ half-edges whose incidence graph on the $k$ vertices is
connected and whose pairs never lie on a single vertex. Define
\[
  G_k^{\mathrm{conn}}(z_1,\ldots,z_k)
  =
  \left(-\frac12\right)^k
  \sum_{\pi\in\mathcal P_k^{\mathrm{conn}}}
  \prod_{\{(i,r),(j,s)\}\in\pi}
  \left(-\frac1{(z_i-z_j)^2}\right).
\]
\end{definition}

\begin{theorem}[Five-point value]
\label{thm:g5-value}
The finite Wick enumeration gives
\[
  G_5^{\mathrm{conn}}(1,2,3,4,5)
  =
  \frac{775}{5184}
  =
  \frac{5^2\cdot31}{2^6\cdot3^4}.
\]
\end{theorem}

\begin{proof}
There are $9!!=945$ perfect matchings of ten labelled half-edges before
the exclusion of self-pairings and disconnected incidence graphs. Each
surviving term in Definition~\ref{def:exact-wick-sum} is a rational
number at $(z_1,\ldots,z_5)=(1,2,3,4,5)$. Summing those rational terms
gives $775/5184$. The local exact arithmetic implementation is
\texttt{compute/lib/virasoro\_m5\_five\_point.py}; the assertion is
checked by
\texttt{compute/tests/test\_virasoro\_m5\_five\_point.py}.
\end{proof}

\begin{remark}[Comparison with $m_5$]
\label{rem:g5-m5-comparison}
The coefficient $m_5=-80/9$ is the Newton coefficient $a_3$ in
Theorem~\ref{thm:c1-shadow-tower-through-eight}. The rational number
$775/5184$ is a connected Wick cumulant at a chosen configuration. A
comparison from the cumulant to the minimal $A_\infty$ operation uses
the Virasoro normalisation, projection to the $T$-line, the level-four
constant $S_4=10/27$, and the Arnold residue map on configuration
space. The finite Wick value is compute-backed independently; the
residue normalisation is the comparison datum.
\end{remark}

\section{$A_\infty$ Coproduct Residues}

Let $B^{\mathrm{ord}}(\mathcal A)$ be the ordered bar complex of the
vertex algebra on the completed $\mathcal W_{1+\infty}$ line. Its
homotopy-transferred coproduct has the form
\[
  \Delta^{A_\infty}
  =
  \Delta^{Y^+}
  +\sum_{k\geq3}\hbar^{k-2}\delta^{(k)}.
\]
The symbol $\Delta^{Y^+}$ denotes the primitive positive-half Hall
coproduct. The completed double is the layer on which the vertex
algebra is read.

\begin{proposition}[Scalar residue normalisation]
\label{prop:scalar-residue-normalisation}
On the $T$-line quotient,
\[
  \delta^{(k)}(T)
  =
  S_k\,\Delta_{\mathrm{sym}}^{(k)}(T)
  +(\text{traceless primary terms}),
\]
where $\Delta_{\mathrm{sym}}^{(k)}$ is the symmetrised $k$-fold
positive-half coproduct. Under the cyclic corolla pairing,
$m_k(T,\ldots,T)=kS_kT$ for $k\geq4$, while
$m_3(T,T,T)=-S_3T$ in the Virasoro associator convention.
\end{proposition}

\begin{proof}
The unit projection on $B^{\mathrm{ord}}(\mathcal A)$ kills the
traceless primary terms and leaves the scalar coefficient of the
symmetrised $k$-leg corolla. For $k\geq4$, the cyclic orbit of the
outgoing leg has length $k$, giving the factor $k$. The ternary
operation is fixed by the Gerstenhaber associator sign convention
$m_3(T,T,T)=-\alpha T$ with $S_3=\alpha$.
\end{proof}

\section{MacMahon and Bar Euler Scope}

The point DT vertex of $\mathbb C^3$ is the MacMahon product
\[
  M(q)=\prod_{n\geq1}(1-q^n)^{-n}.
\]
The ordered bar Euler characteristic of the positive Hall half and the
graded character of $Y^+(\widehat{\mathfrak{gl}}_1)$ are $q$-series
statements. The Virasoro shadow tower is an arity series in $t$ governed
by $\sqrt{Q_1(t)}$. A comparison between these objects requires the
choice of chamber, Behrend sign, grading variable, and projection map.
The equality
\[
  \mathrm{CoHA}(\mathbb C^3)=Y^+(\widehat{\mathfrak{gl}}_1)
\]
lives on the positive associative Hall side; the
$\mathcal W_{1+\infty}$ algebra is the completed double or centre
evaluation.

\section{$K3\times E$ and CHL Constants}

The $c=1$ Virasoro tower above belongs to the $\mathbb C^3$ completion.
The $K3\times E$ Borcherds and CHL constants live in a different
construction. For $Y=K3\times E$,
\[
\begin{array}{c|c|c}
\text{quantity} & \text{value} & \text{source}\\
\hline
\kappa_{\mathrm{cat}}(Y) & 0 &
\chi(\mathcal O_{K3})\chi(\mathcal O_E)=2\cdot0\\
\kappa_{\mathrm{ch}}(Y) & 0 &
\sum_q(-1)^qh^{0,q}(Y)=1-1+1-1\\
\kappa_{\mathrm{ch}}^{\mathrm{Heis}}(Y) & 3 &
\text{K3-fibre Heisenberg specialisation}\\
\kappa_{\mathrm{BKM}}(\Delta_5) & 5 &
c_1(0)/2=10/2\\
\kappa_{\mathrm{fiber}}(Y) & 24 &
\mathrm{rk}\,\widetilde H(K3,\mathbb Z)
\end{array}
\]
These are parallel scalar readings with separate comparison maps.

On the primitive CHL denominator ladder
$N\in\{1,2,3,4,6\}$,
\[
  \bigl(c_1(0),c_2(0),c_3(0),c_4(0),c_6(0)\bigr)
  =(10,8,6,4,2),
\]
so
\[
  \bigl(\kappa_{\mathrm{BKM}}(\Phi_1),
        \kappa_{\mathrm{BKM}}(\Phi_2),
        \kappa_{\mathrm{BKM}}(\Phi_3),
        \kappa_{\mathrm{BKM}}(\Phi_4),
        \kappa_{\mathrm{BKM}}(\Phi_6)\bigr)
  =(5,4,3,2,1).
\]
The law is the Borcherds weight identity
\[
  \kappa_{\mathrm{BKM}}(\Phi_N)=\frac{c_N(0)}2.
\]
The Gritsenko-Clery diagonal-divisor catalogue is a separate
eight-row table with
\[
  c_R(0)=(10,4,6,2,4,1,3,2),\qquad
  \kappa_{\mathrm{BKM}}=(5,2,3,1,2,1/2,3/2,1).
\]

\section{Proof Boundaries}

\begin{enumerate}
\item The values $m_5,m_6,m_7,m_8$ and $S_5,S_6,S_7,S_8$ are proved by
the square-root recurrence from the three Virasoro inputs
$(\kappa_{\mathrm{ch}},\alpha,S_4)$.
\item The value $G_5^{\mathrm{conn}}(1,2,3,4,5)=775/5184$ is an exact
finite Wick enumeration. It is compute-backed independently of the
Newton recurrence.
\item The all-order non-termination of the $c=1$ tower follows from
the fact that $Q_1(t)$ is a non-square quadratic. The displayed
nonzero coefficients are verified only through arity eight.
\item The hCS Feynman reading is conditional on the chosen
Costello-Gaiotto boundary-algebra identification and propagator
normalisation. The rational constants above are independent of that
physical interpretation.
\item The MacMahon product, ordered bar Euler characteristic, Virasoro
shadow tower, and CHL Borcherds constants occupy different graded
objects. Their comparison maps must be named before numerical values
are identified.
\end{enumerate}

\end{document}
